\section{Scheduling and Schedulability}

\subsection{Periodic tasks}
\begin{frame}{Periodic tasks and processor demand}
  Task $i$: period $T_i$, execution time $C_i$, implicit deadline $D_i = T_i$.
  \[ U_i = \frac{C_i}{T_i}, \qquad U = \sum_i U_i \]
  \begin{itemize}
    \item Example: $T_1 = 200$\,ms, $C_1 = 100$\,ms $\Rightarrow$ 50\,\%; $T_2 = 100$\,ms, $C_2 = 50$\,ms $\Rightarrow$ 50\,\%.
    \item Total 100\,\% is the theoretical maximum: no slack, no room for interrupts or overhead.
  \end{itemize}
\end{frame}

\subsection{Rate monotonic}
\begin{frame}{Rate Monotonic (RM)}
  \begin{itemize}
    \item Fixed priorities: the shorter the period, the higher the priority.
    \item Sufficient schedulability test (Liu \& Layland):
    \[ U \le n\left(2^{1/n} - 1\right) \xrightarrow{n\to\infty} \ln 2 \approx 69.3\,\% \]
    \item Exact test: worst-case response time
    \[ R_i = C_i + \sum_{j \in hp(i)} \left\lceil \frac{R_i}{T_j} \right\rceil C_j \]
    (solve iteratively, schedulable if $R_i \le D_i$).
  \end{itemize}
\end{frame}

\subsection{Earliest deadline first}
\begin{frame}{Earliest Deadline First (EDF)}
  \begin{itemize}
    \item Dynamic priorities: the job with the nearest absolute deadline runs first.
    \item For implicit deadlines: schedulable if and only if $U \le 1$.
    \item Not directly supported by FreeRTOS (fixed priorities), but the FreeRTOS priority scheme maps naturally to RM.
  \end{itemize}
\end{frame}

\scriptonly{%
Exercise: implement the two-task example above on the Arduino, assign priorities according to RM, and verify the response times with the logic analyzer.
}
